% Описание текущих реализаций smart_weights и ADP/engine/box_kernel.py.
% Файл использует обозначения и макросы из tex/manifold_v2.tex
% и должен подключаться после него.

\section{Smart weights и box/plateau-ядра}
\label{sec:smart-weights-box-kernel}

В коде ADP реализованы два независимых способа ускорить построение
локальных весов:
\begin{enumerate}
\item режим \texttt{smart\_weights} сохраняет ядро Эпанечникова и
      вычисляет в точности те же плотные блоки весов, но старается не
      считать заведомо нулевые элементы;
\item модуль \texttt{ADP/engine/box\_kernel.py} задаёт отдельные
      box- и plateau-ядра и строит разреженные списки соседей.
\end{enumerate}
Эти режимы не складываются друг с другом. Конфигурация запрещает
\texttt{smart\_weights=True} для box- и plateau-ядер. Обе ветви
сходятся только на входе батчевого вычисления статистик из
раздела~\ref{sec:adp-batched-statistics}.

\subsection{Общая геометрия локализации}

Для наблюдения $\Xv_i$ и центра $\xv_j$ положим
\begin{equation}
    \Delta_{ij}=\Xv_i-\xv_j,
    \qquad
    D_{ij}^2=\|\Delta_{ij}\|^2.
\label{eq:smart-delta-distance}
\end{equation}

В single-index модели текущий индекс $\betav$ имеет единичную норму.
При bandwidth $\hiso$ и параметре анизотропии $\aniso$ аргумент ядра
равен
\begin{equation}
    q_{ij}^{\mathrm{SI}}
    =
    \frac{
        \aniso^2D_{ij}^2
        +\langle\Delta_{ij},\betav\rangle^2
    }{\hiso^2}.
\label{eq:smart-single-q}
\end{equation}
В Python параметрам $\hiso$ и $\aniso$ соответствуют соответственно
\texttt{h} и \texttt{rho}.

Для multi-index модели используем обозначения раздела~\ref{SEDRsa}:
\begin{equation}
    \PEDR\PEDR^{\mathsf T}=\Id_{\dimm},
    \qquad
    \Lambda=\operatorname{diag}
    (\lambda_1,\ldots,\lambda_{\dimm}).
\end{equation}
В коде ортонормированный базис хранится столбцами, то есть массив
\texttt{basis} соответствует $\PEDR^{\mathsf T}$. Введём
\begin{align}
    R_{ij}^2
    &=\|\PEDR\Delta_{ij}\|^2,
    \\
    L_{ij}^2
    &=\|\Lambda^{1/2}\PEDR\Delta_{ij}\|^2,
    \\
    O_{ij}^2
    &=\max\{D_{ij}^2-R_{ij}^2,0\}.
\end{align}
Последний максимум устраняет только малые отрицательные значения,
возникающие из-за округления. Аргумент ядра равен
\begin{equation}
    q_{ij}^{\mathrm{MI}}
    =
    \frac{\aniso^2O_{ij}^2+L_{ij}^2}{\hiso^2}.
\label{eq:smart-multi-q}
\end{equation}
В multi-index коде используются имена
\begin{equation}
    \aniso\longleftrightarrow\texttt{alpha},
    \qquad
    \operatorname{diag}(\Lambda)
    \longleftrightarrow\texttt{eigenvalues}.
\end{equation}

Во всех режимах центры обрабатываются батчами
$\mathcal B\subset\JJJ$ размера не более \texttt{batch\_size}.
Выход каждого weight-builder имеет начальный индекс батча и данные
для его центров; полный массив временных весов
$\nJJJ\times n$ внутри итерации не создаётся.

\subsection{Режим \texttt{smart\_weights}}

Этот режим разрешён только для реализованного в
\texttt{ADP\_Config.py} ядра Эпанечникова
\begin{equation}
    \Kern_{\mathrm E}(q)=(1-q^2)_{+}.
\label{eq:smart-epanechnikov}
\end{equation}
Следовательно, вес положителен ровно при $q<1$. Умный режим меняет
только порядок вычислений, но не формулу $\weight_{ij}$.

\subsubsection{Single-index screening}

Обозначим
\begin{equation}
    P_{ij}^2=\langle\Delta_{ij},\betav\rangle^2.
\end{equation}
Так как $\|\betav\|=1$, имеем $D_{ij}^2\geq P_{ij}^2$, поэтому
из \eqref{eq:smart-single-q} следует точная нижняя оценка
\begin{equation}
    q_{ij}^{\mathrm{SI}}
    \geq
    \frac{(1+\aniso^2)P_{ij}^2}{\hiso^2}.
\end{equation}
Ненулевой вес возможен только для маски
\begin{equation}
    (1+\aniso^2)P_{ij}^2<\hiso^2.
\label{eq:smart-single-mask}
\end{equation}
Сначала код вычисляет проекции всех наблюдений и центров на
$\betav$, а затем строит маску \eqref{eq:smart-single-mask} для
текущего батча. Из-за cancellation в формуле Грама вычисленное
$D_{ij}^2$ иногда может оказаться меньше $P_{ij}^2$. Такие пары
принудительно добавляются в маску, чтобы численная ошибка не создала
ложный нулевой вес.

Если доля кандидатов не меньше $1/64$, булев индекс и gather обычно
дороже плотных NumPy-операций. Поэтому код возвращается к плотной
формуле
\begin{equation}
    W_{\mathcal B}
    =
    \Kern_{\mathrm E}\!\left(
        \frac{\aniso^2D_{\mathcal B}^2+P_{\mathcal B}^2}
             {\hiso^2}
    \right).
\end{equation}
Только при доле кандидатов строго меньше $1/64$ создаётся нулевой
блок весов, а полный аргумент ядра вычисляется по элементам маски.

\subsubsection{Multi-index screening}

Из неотрицательности $O_{ij}^2$ в
\eqref{eq:smart-multi-q} следует
\begin{equation}
    q_{ij}^{\mathrm{MI}}
    \geq\frac{L_{ij}^2}{\hiso^2}.
\end{equation}
Поэтому точная маска кандидатов имеет вид
\begin{equation}
    L_{ij}^2<\hiso^2.
\label{eq:smart-multi-mask}
\end{equation}
Проекции на столбцы базиса вычисляются один раз, после чего
$R_{ij}^2$ и $L_{ij}^2$ накапливаются по $\dimm$ компонентам.
Дальнейшая развилка та же: при плотности маски не меньше $1/64$
вычисляется полный блок, иначе $O_{ij}^2$ и значение ядра считаются
только для кандидатов.

Тест \texttt{test\_ADP\_smart.py} сравнивает обычную и умную ветви
для single- и multi-index, в том числе на очень разреженной опоре и
при большом общем сдвиге данных. Контракт режима состоит в совпадении
всех весов с плотной ветвью до ошибки округления.

\subsubsection{Точный выбор $\aniso$ для ядра Эпанечникова}

Без \texttt{smart\_weights} новые $\aniso$ в single-index и
multi-index выбираются бинарным поиском. На каждой пробе заново
вычисляется средняя локальная масса
\begin{equation}
    \frac1{\nJJJ}
    \sum_{j\in\JJJ}\sum_{i=1}^n
    \Kern_{\mathrm E}(q_{ij})
    \geq\nmin.
\label{eq:smart-mean-mass}
\end{equation}

Умная ветвь использует кусочно-квадратичную структуру
\eqref{eq:smart-epanechnikov}. Для обеих моделей положим
\begin{equation}
    z=\aniso^2,
    \qquad
    q_{ij}(z)=\frac{a_{ij}z+b_{ij}}{\hiso^2},
\end{equation}
где
\begin{equation}
    (a_{ij},b_{ij})
    =
    \begin{cases}
    (D_{ij}^2,P_{ij}^2), & \text{single-index},\\
    (O_{ij}^2,L_{ij}^2), & \text{multi-index}.
    \end{cases}
\end{equation}
Требуемая суммарная масса равна
$M_{\star}=\nmin\nJJJ$. Для фиксированного набора активных пар
\begin{equation}
    \mathcal A(z)
    =\{(j,i):a_{ij}z+b_{ij}<\hiso^2\}
\end{equation}
масса выражается точно:
\begin{align}
    M(z)
    &=
    \sum_{(j,i)\in\mathcal A(z)}
    \left[1-\frac{(a_{ij}z+b_{ij})^2}{\hiso^4}\right]
    \\
    &=
    |\mathcal A(z)|
    -\frac{
        z^2\sum a_{ij}^2
        +2z\sum a_{ij}b_{ij}
        +\sum b_{ij}^2
    }{\hiso^4}.
\label{eq:smart-piecewise-mass}
\end{align}
Активный набор меняется только в точках
\begin{equation}
    \tau_{ij}=\frac{\hiso^2-b_{ij}}{a_{ij}}.
\label{eq:smart-breakpoints}
\end{equation}
Функция \texttt{select\_optimal\_alpha} отбрасывает пары с
$b_{ij}\geq\hiso^2$, сортирует остальные $\tau_{ij}$ и один раз
строит суффиксные суммы
$\sum a_{ij}^2$, $\sum a_{ij}b_{ij}$ и $\sum b_{ij}^2$.
Затем бинарным поиском по точкам разрыва находится интервал,
содержащий наибольшее допустимое $z$, и внутри него решается одно
квадратное уравнение $M(z)=M_{\star}$. Возвращается
$\aniso=\sqrt z$.

Если $M(1)\geq M_{\star}$, сразу возвращается $1$. Если
$M(0)<M_{\star}$, допустимого масштаба нет и возвращается
\texttt{None}; outer-loop завершится с причиной
\texttt{local\_mass\_limit}.

У этого ускорения есть цена: во время выбора масштаба существуют
массивы размера $\nJJJ\times n$, маска активных пар, отсортированные
breakpoints и суффиксные суммы. Показатель выделений отдельной стадии
нельзя интерпретировать как одновременный пик памяти всего fit.

\subsection{Разреженный backend \texttt{box\_kernel.py}}

\subsubsection{Поддерживаемые ядра}

Box-ядро задаётся формулой
\begin{equation}
    \Kern_{\mathrm{box}}(q)
    =\mathbf 1\{q<1\}.
\label{eq:box-kernel}
\end{equation}
Plateau-ядро имеет параметр $\tau_{K}\in(0,1)$. Для
\begin{equation}
    s=\frac{q-\tau_K}{1-\tau_K}
\end{equation}
оно равно
\begin{equation}
    \Kern_{\mathrm{pl}}(q)
    =
    \begin{cases}
    1, & q\leq\tau_K,\\
    1-10s^3+15s^4-6s^5, & \tau_K<q<1,\\
    0, & q\geq1.
    \end{cases}
\label{eq:plateau-kernel}
\end{equation}
Полиномиальный участок имеет нулевые первую и вторую производные на
обеих границах. Функция \texttt{sparse\_kernel\_parameters}
распознаёт только сам \texttt{box\_kernel}, стандартный
\texttt{plateau\_kernel} с $\tau_K=1/2$ или результат
\texttt{make\_plateau\_kernel($\tau_K$)}.

\subsubsection{Объект \texttt{NeighborhoodEngine}}

Движок получает неизменные $X$, центры и размер блока. При создании он
строит \texttt{scipy.spatial.cKDTree} по наблюдениям. Для box-ядра
движок дополнительно может сохранить полную матрицу
$D^2\in\R^{\nJJJ\times n}$, но только если она занимает не более
$64\,\mathrm{MiB}$. Если лимит превышен или матрица не конечна,
используется разреженный поиск через деревья.

Внешняя граница обоих sparse-ядер одинакова: пара входит в support
ровно при $q_{ij}<1$. Деревья используются только для получения
супермножества кандидатов; окончательное решение всегда принимает
точная формула \eqref{eq:smart-single-q} или
\eqref{eq:smart-multi-q}.

\paragraph{Изотропные окрестности.}
Для начальной локализации выполняется radius query
\begin{equation}
    \|\Delta_{ij}\|<\hiso,
\end{equation}
после чего вычисляется $q_{ij}=D_{ij}^2/\hiso^2$ и оставляются пары
с $q_{ij}<1$. При наличии box-кэша support извлекается сразу из строк
матрицы $D^2$.

\paragraph{Single-index окрестности.}
Из \eqref{eq:smart-single-q} следуют два необходимых условия
\begin{equation}
    |\langle\Delta_{ij},\betav\rangle|
    <\frac{\hiso}{\sqrt{1+\aniso^2}},
    \qquad
    \|\Delta_{ij}\|<\frac{\hiso}{\aniso}
    \quad(\aniso>0).
\label{eq:box-single-screens}
\end{equation}
Первая выборка строится бинарным поиском по отсортированным проекциям
$\betav^{\mathsf T}X$. Вторая получается запросом к исходному
\texttt{cKDTree}. Пересечение списков проверяется точной формулой
$q_{ij}^{\mathrm{SI}}<1$. При $\aniso=0$ евклидово ограничение
отсутствует. Для box-ядра с доступным кэшем $D^2$ support строится
непосредственно без этих деревьев.

\paragraph{Multi-index окрестности.}
Сначала строятся преобразованные координаты
\begin{equation}
    Z_i=\Lambda^{1/2}\PEDR\Xv_i,
    \qquad
    Z_j^c=\Lambda^{1/2}\PEDR\xv_j.
\end{equation}
Условие $L_{ij}^2<\hiso^2$ превращается в radius query радиуса
$\hiso$ в $\dimm$-мерном пространстве $Z$. Дополнительное евклидово
ограничение имеет радиус
\begin{equation}
    \frac{\hiso}{
        \sqrt{\min\{\aniso^2,\lambda_{\min}\}}
    }
\end{equation}
и применяется, только если знаменатель положителен. Пересечение
кандидатов снова фильтруется точным условием
$q_{ij}^{\mathrm{MI}}<1$. Для box-кэша компоненты
$O_{ij}^2$ и $L_{ij}^2$ вычисляются матрично по всем парам.

\subsubsection{Формат \texttt{SparseNeighborhoodBlock}}

Найденные окрестности делятся на блоки центров. Каждый неизменяемый
объект хранит
\begin{itemize}
\item \texttt{start} --- индекс первого центра;
\item \texttt{indptr} и \texttt{indices} --- границы строк и индексы
      всех пар с $q_{ij}<1$;
\item \texttt{boundary\_positions} и\\
      \texttt{boundary\_weights} --- только plateau-пары с
      $\tau_K<q_{ij}<1$;
\item режим ядра, $\tau_K$ и общее число наблюдений.
\end{itemize}
Для box-ядра все сохранённые веса неявно равны единице, поэтому
массивы границы пусты. Для plateau-ядра пары с $q_{ij}\leq\tau_K$
также имеют неявный вес один; явно хранятся только значения на
полиномиальной границе.

Масса строки получается из числа индексов с поправкой
\begin{equation}
    \mu_j
    =|\mathcal N_j|
     +\sum_{i\in\mathcal B_j}(\weight_{ij}-1),
\end{equation}
где $\mathcal B_j$ --- boundary layer plateau-ядра. Метод
\texttt{row} восстанавливает одну строку, а \texttt{padded}
переводит блок в прямоугольные массивы индексов и весов для
батчевого ядра статистик.

Метод \texttt{cache\_key} кодирует точный support строки, а для
plateau также позиции и значения граничных весов. Эти ключи нужны
для диагностики повторного использования. Текущий
\texttt{SparseStatisticsCache} считает совпавшие строки между двумя
вызовами, но не хранит и не переиспользует численные статистики.

\subsubsection{Выбор начального bandwidth}

Для box-ядра условие достаточной локализации является построчным:
\begin{equation}
    \#\{i:q_{ij}<1\}\geq\lceil\nmin\rceil
\end{equation}
должно выполняться хотя бы для
\begin{equation}
    \left\lceil0.9\,\nJJJ\right\rceil
\end{equation}
центров. При доступном кэше $D^2$ функция
\texttt{search\_sparse\_bandwidth} выбирает соответствующую
порядковую статистику сначала в каждой строке, затем по центрам.
Поскольку support использует строгое неравенство, найденный радиус
сдвигается через \texttt{nextafter} в сторону $+\infty$.

Если кэш превысил $64\,\mathrm{MiB}$, тот же критерий проверяется
построением sparse-блоков и бинарным поиском по $\hiso$. Для
plateau-ядра используется другой критерий --- средняя взвешенная
масса по всем центрам не меньше $\nmin$ --- и бинарный поиск.

\subsubsection{Выбор $\aniso$ для box и plateau}

Для box-ядра с доступным $D^2$ используются функции
\begin{center}
    \texttt{select\_single\_box\_scale},\\
    \texttt{select\_multi\_box\_scale}.
\end{center}
Они используют точные пороги
\begin{align}
    \aniso^2
    &<\frac{\hiso^2-P_{ij}^2}{D_{ij}^2},
    &&\text{single-index},
    \\
    \aniso^2
    &<\frac{\hiso^2-L_{ij}^2}{O_{ij}^2},
    &&\text{multi-index}.
\end{align}
Из порогов выбирается наибольший масштаб, при котором не менее $90\%$
центров имеют хотя бы $\lceil\nmin\rceil$ соседей. Значение
сдвигается через \texttt{nextafter} к нулю, чтобы удовлетворить
строгому условию $q<1$.

Multi-index selector одновременно сохраняет найденный support.
Следующий вызов \texttt{multi\_blocks} с теми же байтами базиса,
собственными значениями, $\hiso$ и $\aniso$ использует этот список
без повторного вычисления компонент расстояния.

Если полный кэш недоступен, selector возвращает
\texttt{NotImplemented}; модель переходит к
\texttt{search\_sparse\_scale}. Этот поиск проверяет точки $0$ и $1$,
а затем выполняет бинарный поиск с точностью
$\sqrt{\varepsilon_{\mathrm{mach}}}$, каждый раз заново строя
sparse-блоки. Для box сохраняется построчный критерий $90\%$, для
plateau используется средняя взвешенная масса. Невыполнимость уже
при $\aniso=0$ возвращает \texttt{None}.

\subsection{Подключение к outer-loop}

При обычном ядре модель заранее вычисляет $D^2$ и вызывает одну из
функций
\begin{center}
    \texttt{calculate\_weight},\qquad
    \texttt{calculate\_multi\_weight}
\end{center}
с флагом \texttt{smart}. Эти функции выдают плотные блоки
$B\times n$, которые сразу потребляются
\texttt{calculate\_statistics}.

При box/plateau модель не создаёт обычную матрицу $D^2$ и вместо неё
создаёт один \texttt{NeighborhoodEngine}. Тот же движок используется
для sparse local-linear initialization, выбора начального $\hiso$,
построения окрестностей на каждой итерации и выбора следующего
$\aniso$. Объекты
\begin{center}
    \texttt{SparseNeighborhoodBlock}
\end{center}
передаются в \texttt{calculate\_statistics}, где превращаются в
padded-батчи.

После решения локальной задачи модель обновляет
\begin{equation}
    \hiso_{k+1}=\hiso_k/a
\end{equation}
и выбирает наибольший допустимый $\aniso_{k+1}$. Если допустимого
значения нет, итерации завершаются с причиной
\begin{center}
    \texttt{stop\_reason=local\_mass\_limit}.
\end{center}
В trace для sparse-ветви
записываются число рёбер support, число plateau boundary-рёбер и
число строк с тем же точным support, что на предыдущем вызове.

Построение окрестностей остаётся CPU-операцией и в GPU-режиме.
На CUDA переносятся padded-блоки и дальнейшая линейная алгебра
статистик. Автоматического перехода между smart- и sparse-ветвями
нет: конкретный путь полностью определяется \texttt{ADP\_Config}.
